\documentclass[10pt,a4paper]{amsart}
\usepackage{amsmath,amssymb,amsthm,graphicx}
\usepackage{verbatim}
\newenvironment{reference}{\comment}{\endcomment}
%\newenvironment{reference}{}{}
\newenvironment{slogan}{\comment}{\endcomment}
\newenvironment{history}{\comment}{\endcomment}

% For commutative diagrams we use Xy-pic
\usepackage[all]{xy}

% We use 2cell for 2-commutative diagrams.
\xyoption{2cell}
% Theorem environments.
%
\theoremstyle{plain}
\newtheorem{theorem}[subsection]{Theorem}
\newtheorem{proposition}[subsection]{Proposition}
\newtheorem{lemma}[subsection]{Lemma}

\theoremstyle{definition}
\newtheorem{definition}[subsection]{Definition}
\newtheorem{example}[subsection]{Example}
\newtheorem{exercise}[subsection]{Exercise}
\newtheorem{situation}[subsection]{Situation}

\theoremstyle{remark}
\newtheorem{remark}[subsection]{Remark}
\newtheorem{remarks}[subsection]{Remarks}

\numberwithin{equation}{subsection}

% Macros
%
\def\lim{\mathop{\mathrm{lim}}\nolimits}
\def\colim{\mathop{\mathrm{colim}}\nolimits}
\def\Spec{\mathop{\mathrm{Spec}}}
\def\Hom{\mathop{\mathrm{Hom}}\nolimits}
\def\Ext{\mathop{\mathrm{Ext}}\nolimits}
\def\SheafHom{\mathop{\mathcal{H}\!\mathit{om}}\nolimits}
\def\SheafExt{\mathop{\mathcal{E}\!\mathit{xt}}\nolimits}
\def\Sch{\mathit{Sch}}
\def\Mor{\mathop{\mathrm{Mor}}\nolimits}
\def\Ob{\mathop{\mathrm{Ob}}\nolimits}
\def\Sh{\mathop{\mathit{Sh}}\nolimits}
\def\NL{\mathop{N\!L}\nolimits}
\def\CH{\mathop{\mathrm{CH}}\nolimits}
\def\proetale{{pro\text{-}\acute{e}tale}}
\def\etale{{\acute{e}tale}}
\def\QCoh{\mathit{QCoh}}
\def\Ker{\mathop{\mathrm{Ker}}}
\def\Im{\mathop{\mathrm{Im}}}
\def\Coker{\mathop{\mathrm{Coker}}}
\def\Coim{\mathop{\mathrm{Coim}}}

% Boxtimes
%
\DeclareMathSymbol{\boxtimes}{\mathbin}{AMSa}{"02}

%
% Macros for moduli stacks/spaces
%
\def\QCohstack{\mathcal{QC}\!\mathit{oh}}
\def\Cohstack{\mathcal{C}\!\mathit{oh}}
\def\Spacesstack{\mathcal{S}\!\mathit{paces}}
\def\Quotfunctor{\mathrm{Quot}}
\def\Hilbfunctor{\mathrm{Hilb}}
\def\Curvesstack{\mathcal{C}\!\mathit{urves}}
\def\Polarizedstack{\mathcal{P}\!\mathit{olarized}}
\def\Complexesstack{\mathcal{C}\!\mathit{omplexes}}
% \Pic is the operator that assigns to X its picard group, usage \Pic(X)
% \Picardstack_{X/B} denotes the Picard stack of X over B
% \Picardfunctor_{X/B} denotes the Picard functor of X over B
\def\Pic{\mathop{\mathrm{Pic}}\nolimits}
\def\Picardstack{\mathcal{P}\!\mathit{ic}}
\def\Picardfunctor{\mathrm{Pic}}
\def\Deformationcategory{\mathcal{D}\!\mathit{ef}}
\usepackage{cleveref}
\crefname{originallemma}{original lemma}{original lemmas}
\Crefname{originallemma}{Original lemma}{Original lemmas}
\makeatletter
\newlabel{lemma-point-site-topos}{{00YA}{source}{}{sourcepointer-1}{}}
\newlabel{lemma-point-site-topos@cref}{{[originallemma][0][]00YA}{[1][0][]source}{}{}{}}
\makeatother
\crefname{originallemma}{original lemma}{original lemmas}
\Crefname{originallemma}{Original lemma}{Original lemmas}
\makeatletter
\newlabel{lemma-technical-pu}{{00XM}{source}{}{sourcepointer-2}{}}
\newlabel{lemma-technical-pu@cref}{{[originallemma][0][]00XM}{[1][0][]source}{}{}{}}
\makeatother
\crefname{originallemma}{original lemma}{original lemmas}
\Crefname{originallemma}{Original lemma}{Original lemmas}
\makeatletter
\newlabel{lemma-sheafification-exact}{{00WJ}{source}{}{sourcepointer-3}{}}
\newlabel{lemma-sheafification-exact@cref}{{[originallemma][0][]00WJ}{[1][0][]source}{}{}{}}
\makeatother
\title{Points of a localized site lying over a given point}
\author{The Stacks Project Authors}
\date{stacks-project:sites.tex:local-snapshot-20261012; GNU FDL 1.2 or later}
\begin{document}
\maketitle
\noindent\textit{Editorial note.} All delivered mathematical text is the author's own, copied from the pinned Stacks Project chapter file sites-cohomology.tex; no statement or proof step was rewritten, reordered or removed.\par\medskip
\noindent\textit{Editorial note.} The proof is the author's complete argument; every step the author writes out is inside the counted span and no closing sentence is a placeholder.\par\medskip
\noindent\textit{Editorial note.} Every labelled reference occurring inside the counted spans is carried in assembly.external\_references as a non-load-bearing navigation pointer (same-chapter ones at their definition site, cross-chapter ones at the invocation inside the counted span), so no load-bearing mathematics is replaced by a stub. No bibliography citation occurs inside the counted spans.\par\medskip
\noindent\textit{Editorial note.} The checked chapter list for same-name collisions is the set of stacks chapters staged in this repository: chow, cohomology, divisors, etale-cohomology, flat, more-algebra, more-morphisms, morphisms, spaces-more-morphisms, varieties together with this card's chapter sites-cohomology.\par\medskip
\noindent\textit{Editorial note.} The preamble is the chapter's own retained minimal preamble (sources/stacks-sites-cohomology/retained-preamble.tex), a verbatim excerpt of the native preamble: the Xy-pic setup and the native theorem declarations with the mathematical macros.\par\medskip
\noindent\textit{Allowed standard prerequisites (editorial).}\begin{itemize}
\item The point of a localized site attached to an element. Same chapter, Lemma lemma-point-localize, recorded as the context of this card. Supplies the construction whose fibres are described by the correspondence.
\item Points of a site versus points of its topos. Same chapter, Lemma lemma-point-site-topos. Identifies the values of the functor attached to a point with the pullback of representable sheaves, which is what is used to compute the diagrams.
\end{itemize}\medskip
\ifdefined\hypertarget\hypertarget{sourcepointer-1}{}\fi
\noindent\textit{Source reference.} Original lemma 00YA: sources/stacks-sites/pinned-original-source.tex, source lines 8981--8981. Navigation pointer invoked at proof line 8981; the counted span contains the pointer only and the cited result is not reproduced inside it, so the pointer is registered as a non-load-bearing navigation stub in the established form for Stacks material. No bibliography citation is involved.\par\medskip
\ifdefined\hypertarget\hypertarget{sourcepointer-2}{}\fi
\noindent\textit{Source reference.} Original lemma 00XM: sources/stacks-sites/pinned-original-source.tex, source lines 9006--9006. Navigation pointer invoked at proof line 9006; the counted span contains the pointer only and the cited result is not reproduced inside it, so the pointer is registered as a non-load-bearing navigation stub in the established form for Stacks material. No bibliography citation is involved.\par\medskip
\ifdefined\hypertarget\hypertarget{sourcepointer-3}{}\fi
\noindent\textit{Source reference.} Original lemma 00WJ: sources/stacks-sites/pinned-original-source.tex, source lines 9009--9009. Navigation pointer invoked at proof line 9009; the counted span contains the pointer only and the cited result is not reproduced inside it, so the pointer is registered as a non-load-bearing navigation stub in the established form for Stacks material. No bibliography citation is involved.\par\medskip
\begin{lemma}
\label{lemma-point-localize}
Let $\mathcal{C}$ be a site. Let $p$ be a point of $\mathcal{C}$ given by
$u : \mathcal{C} \to \textit{Sets}$. Let $U$ be an object of $\mathcal{C}$
and let $x \in u(U)$. The functor
$$
v : \mathcal{C}/U \longrightarrow \textit{Sets}, \quad
(\varphi : V \to U) \longmapsto \{y \in u(V) \mid u(\varphi)(y) = x\}
$$
defines a point $q$ of the site $\mathcal{C}/U$ such that the diagram
$$
\xymatrix{
& \Sh(pt) \ar[d]^p \ar[ld]_q \\
\Sh(\mathcal{C}/U) \ar[r]^{j_U} &
\Sh(\mathcal{C})
}
$$
commutes. In other words
$\mathcal{F}_p = (j_U^{-1}\mathcal{F})_q$ for any
sheaf on $\mathcal{C}$.
\end{lemma}

\begin{lemma}
\label{lemma-points-above-point}
Let $\mathcal{C}$, $p$, $u$, $U$ be as in
Lemma \ref{lemma-point-localize}.
The construction of
Lemma \ref{lemma-point-localize}
gives a one to one correspondence between points $q$
of $\mathcal{C}/U$ lying over $p$ and elements $x$ of $u(U)$.
\end{lemma}

\begin{proof}
Let $q$ be a point of $\mathcal{C}/U$ given by the functor
$v : \mathcal{C}/U \to \textit{Sets}$ such that $j_U \circ q = p$
as morphisms of topoi. Recall that $u(V) = p^{-1}(h_V^\#)$ for any
object $V$ of $\mathcal{C}$, see
Lemma \ref{lemma-point-site-topos}.
Similarly $v(V/U) = q^{-1}(h_{V/U}^\#)$ for any object
$V/U$ of $\mathcal{C}/U$.
Consider the following two diagrams
$$
\vcenter{
\xymatrix{
\Mor_{\mathcal{C}/U}(W/U, V/U) \ar[r] \ar[d] &
\Mor_\mathcal{C}(W, V) \ar[d] \\
\Mor_{\mathcal{C}/U}(W/U, U/U) \ar[r] &
\Mor_\mathcal{C}(W, U)
}
}
\quad
\vcenter{
\xymatrix{
h_{V/U}^\# \ar[r] \ar[d] & j_U^{-1}(h_V^\#) \ar[d] \\
h_{U/U}^\# \ar[r] & j_U^{-1}(h_U^\#)
}
}
$$
The right hand diagram is the sheafification of the diagram of
presheaves on $\mathcal{C}/U$ which maps $W/U$ to the left hand
diagram of sets. (There is a small technical point to make here, namely, that
we have $(j_U^{-1}h_V)^\# = j_U^{-1}(h_V^\#)$ and similarly for $h_U$, see
Lemma \ref{lemma-technical-pu}.)
Note that the left hand diagram of sets is cartesian.
Since sheafification is exact
(Lemma \ref{lemma-sheafification-exact})
we conclude that the right hand diagram is cartesian.

\medskip\noindent
Apply the exact functor $q^{-1}$ to the right
hand diagram to get a cartesian diagram
$$
\xymatrix{
v(V/U) \ar[r] \ar[d] & u(V) \ar[d] \\
v(U/U) \ar[r] & u(U)
}
$$
of sets. Here we have used that
$q^{-1} \circ j^{-1} = p^{-1}$. Since $U/U$ is a final object of
$\mathcal{C}/U$ we see that $v(U/U)$ is a singleton. Hence the
image of $v(U/U)$ in $u(U)$ is an element $x$, and the top horizontal
map gives a bijection
$v(V/U) \to  \{y \in u(V) \mid y \mapsto x\text{ in }u(U)\}$
as desired.
\end{proof}

\end{document}
